\section{Integração Endo-Prótese e Arcabouços de Modelagem Unificada}

A simbiose profunda entre endodontia e prótese no contexto de um mesmo modelo computacional é uma das perspectivas mais promissoras da odontologia digital, visto que a fragmentação desses saberes na prática clínica rotineira é responsável por grande parte dos insucessos estruturais. Elementos dentários tratados endodonticamente que necessitam de retentores intrarradiculares e cobertura coronária total formam o epicentro desse paradigma integrativo.

O gêmeo digital unificado converte as barreiras informacionais entre as especialidades em metadados coesos. O modelo assimila a geometria cônica do preparo endodôntico (derivada de dados tomográficos segmentados) e a topologia externa recriada virtualmente \cite{alshadidi2024surface}. As interfaces adesivas -- dentina intrarradicular, cimento resinoso, núcleo de preenchimento e substrato do pino (fibra de vidro ou zircônia) -- são todas parametrizadas com coeficientes elásticos e limites de escoamento específicos.

\destaque{A identificação precoce do risco de fratura radicular vertical, evento tipicamente terminal para o órgão dental, é a maior conquista clínica proporcionada pela simulação do gêmeo digital endo-protético.}

Através da modelagem multifísica, simula-se a dissipação de cargas mastigatórias sobre o elemento restaurado, correlacionando o módulo de elasticidade do pino à quantidade volumétrica de tecido dentinário remanescente. O desgaste da dentina pericervical decorrente do acesso pulpar compromete o chamado "efeito férula", crítico para a resistência mecânica. O sistema não apenas avalia a viabilidade do dente, mas sugere ativamente alterações geométricas no preparo protético, visando transferir o pico de tensão da raiz para a interface restauradora, atuando sob o conceito \textit{fail-safe} \cite{xing2024clinical}.

A plataforma OpenCode (detalhada no Capítulo~\ref{ch:opencode}) introduz uma camada arquitetural indispensável para essa orquestração. Baseada na norma FMI (\textit{Functional Mock-up Interface}) e em uma topologia distribuída, o gêmeo digital do dente trafega por sistemas de agentes virtuais heterogêneos. O agente especializado em anatomia endodôntica submete os volumes tomográficos ao agente de design protético CAD, gerando ciclos automáticos de otimização estrutural até a obtenção de um plano de tratamento único e metodologicamente rigoroso \cite{robles2023opentwins}.
